% ========================
% Typography and layout
% ========================
\usepackage[margin=1in]{geometry}
\usepackage{microtype}
\usepackage{ragged2e}
\linespread{1.05}
\setlength{\parindent}{0pt}
\setlength{\parskip}{0.45em}

% ========================
% Math
% ========================
\usepackage{amsmath,amssymb,amsthm,mathtools,amsfonts,bbm,mathrsfs}
\usepackage{bm}

% ========================
% Figures, tables, and lists
% ========================
\usepackage{graphicx}
\usepackage{booktabs}
\usepackage{array}
\usepackage{longtable}
\usepackage{enumitem}
\setlist{nosep}

% ========================
% Citations
% ========================
\usepackage[numbers,sort&compress]{natbib}
\setlength{\bibhang}{1.5em}
\setlength{\bibsep}{0pt}
\let\oldthebibliography\thebibliography
\let\endoldthebibliography\endthebibliography
\renewenvironment{thebibliography}[1]
  {\oldthebibliography{#1}%
   \RaggedRight
   \frenchspacing
   \hbadness=10000
   \def\newblock{\unskip\space}%
   \setlength{\parskip}{0pt}%
   \setlength{\itemsep}{0pt}%
   \setlength{\parsep}{0pt}}
  {\endoldthebibliography}

% ========================
% Math operators
% ========================
\DeclareMathOperator{\diag}{diag}
\DeclareMathOperator{\tr}{tr}
\DeclareMathOperator{\rank}{rank}
\DeclareMathOperator{\im}{im}
\DeclareMathOperator{\covop}{Cov}
\DeclareMathOperator{\Covop}{COV}
\DeclareMathOperator{\cum}{Cum}
\DeclareMathOperator{\Cumop}{CUM}

\DeclareMathOperator{\corrop}{Corr}
\DeclareMathOperator{\varop}{Var}
\DeclareMathOperator{\Exp}{Exp}
\DeclareMathOperator{\MSE}{MSE}
\DeclareMathOperator{\MAE}{MAE}
\DeclareMathOperator{\ReLU}{ReLU}
\DeclareMathOperator{\softmax}{\widetilde{\max}}
\DeclareMathOperator*{\argmax}{arg\,max}
\DeclareMathOperator*{\argmin}{arg\,min}

\newcommand{\sphere}{\mathbb S}

% ========================
% Number systems and common objects
% ========================
\newcommand{\A}{\mathbb{A}}
\newcommand{\E}{\mathbb{E}}
\newcommand{\N}{\mathbb{N}}
\newcommand{\Z}{\mathbb{Z}}
\newcommand{\Q}{\mathbb{Q}}
\newcommand{\R}{\mathbb{R}}
\newcommand{\C}{\mathbb{C}}
\newcommand{\PP}{\mathbb{P}}
\newcommand{\sF}{\mathscr F}

\newcommand{\cD}{\mathcal{D}}
\newcommand{\cF}{\mathcal{F}}
\newcommand{\cG}{\mathcal{G}}
\newcommand{\cI}{\mathcal{I}}
\newcommand{\cL}{\mathcal{L}}
\newcommand{\cX}{\mathcal{X}}
\newcommand{\cY}{\mathcal{Y}}
\newcommand{\K}{\kappa}
\newcommand{\M}{\mathbb M}

% ========================
% Calculus, probability, and statistics
% ========================
\newcommand{\dd}{\,\mathrm{d}}
\newcommand{\pd}[2]{\frac{\partial #1}{\partial #2}}
\newcommand{\od}[2]{\frac{\mathrm{d} #1}{\mathrm{d} #2}}
\newcommand{\1}{\mathbbm 1}
\newcommand{\var}[1]{\varop\!\left(#1\right)}
\newcommand{\cov}[2]{\covop\!\left(#1,#2\right)}
\newcommand{\COV}[1]{\Covop\!\left(#1\right)}
\newcommand{\CUM}[1]{\Cumop\!\left(#1\right)}
\newcommand{\corr}[2]{\corrop\!\left(#1,#2\right)}
\newcommand{\pro}[1]{\PP\!\left(#1\right)}

\newcommand{\pt}{\partial}

% ========================
% Paired delimiters
% ========================
\DeclarePairedDelimiter{\norm}{\lVert}{\rVert}
\DeclarePairedDelimiter{\abs}{\lvert}{\rvert}
\DeclarePairedDelimiter{\inner}{\langle}{\rangle}
\DeclarePairedDelimiter{\paren}{\lparen}{\rparen}
\DeclarePairedDelimiter{\bracket}{\lbrack}{\rbrack}
\DeclarePairedDelimiter{\set}{\lbrace}{\rbrace}

\newcommand{\nm}[1]{\norm*{#1}}
\newcommand{\inn}[1]{\inner*{#1}}
\newcommand{\ab}[1]{\abs*{#1}}
\newcommand{\st}[1]{\set*{#1}}
\newcommand{\bk}[1]{\bracket*{#1}}
\newcommand{\pa}[1]{{\paren*{#1}}}
\newcommand{\tx}[1]{\text{#1}}

% Compatibility aliases used in older notes.
\newcommand{\rl}{\ReLU}
\newcommand{\x}{\times}
\newcommand{\ox}{\otimes}
\newcommand{\xx}{\mathcal{X}}
\newcommand{\odx}{\odot}
\newcommand{\dx}{\cdot}
\DeclareMathOperator{\sign}{sign}
\DeclareMathOperator{\mom}{momentum}


\DeclareMathOperator{\vol}{volatility}
\DeclareMathOperator{\dv}{dollar\_volume}
\DeclareMathOperator{\turnover}{turnover}
\DeclareMathOperator{\mf}{market\_factor}
\DeclareMathOperator{\mb}{market\_beta}
\DeclareMathOperator{\sector}{sector}


% ========================
% Theorem-like environments
% ========================
\theoremstyle{plain}
\newtheorem{theorem}{Theorem}[section]
\newtheorem{proposition}[theorem]{Proposition}
\newtheorem{lemma}[theorem]{Lemma}
\newtheorem{corollary}[theorem]{Corollary}

\theoremstyle{definition}
\newtheorem{definition}[theorem]{Definition}
\newtheorem{example}[theorem]{Example}
\newtheorem{assumption}[theorem]{Assumption}

\theoremstyle{remark}
\newtheorem{remark}[theorem]{Remark}
\newtheorem{notation}[theorem]{Notation}

% ========================
% Problem set environments
% ========================
\newcounter{problem}

\newenvironment{problem}[1][]{
  \refstepcounter{problem}
  \par\medskip
  \noindent\textbf{Problem \theproblem. #1}
  \par
}{
  \medskip
}

\newenvironment{solution}[1][]{
  \par\medskip
  \noindent\textbf{Solution #1.}
  \par
}{
  \medskip
}

% ========================
% Links and references
% ========================
\usepackage[hidelinks]{hyperref}
\usepackage[capitalize,nameinlink]{cleveref}
